\chapter[Project Management]{Project Management}{Project Management}\label{projm}


\rule{0.49\textwidth}{0.75pt} $_{\bigcirc}$ \rule{0.49\textwidth}{0.75pt}\\





\section{Feasibility Study}

\subsection{Technical Feasibility}
DAVS has a high technical feasibility because it is composed of well-established building blocks from federated learning, distributed systems, and blockchain-based audit logging. The federated training pipeline can be implemented using standard deep learning frameworks, with hospitals acting as clients that perform local training and transmit model updates to a coordinator. Random Projection-based gradient sketching is computationally lightweight and relies on efficient linear algebra operations while providing substantial communication savings.

The Data-Aware Verifier Selection (DAVS) stage operates on low-dimensional sketches and uses cosine similarity and magnitude consistency to select a representative committee each round. This selection is amnesic (round-local) and avoids the complexity and vulnerability of maintaining long-term reputation histories. The PBFT-inspired validation layer leverages known Byzantine fault-tolerant consensus principles with deterministic update checks and quorum voting. Finally, a permissioned blockchain audit layer can be implemented as a lightweight ledger that stores round metadata and hashes (rather than bulky artifacts), making the solution practical for institutional deployments.
\subsection{Operational Feasibility}

The operational feasibility of DAVS is strong in the context of collaborative medical AI. Hospitals can participate as federated clients, training on-premise data while transmitting only compressed sketches and receiving global model updates. Communication efficiency improves round latency and makes multi-site training more deployable in bandwidth-constrained settings.

From a security and governance perspective, DAVS improves operational reliability by excluding abnormal updates before consensus and by validating global updates through committee voting. The permissioned blockchain audit layer provides immutable provenance: per-round committee selection, PBFT votes, model hashes, and metrics can be audited for forensic investigation and compliance-driven reporting, addressing a major gap in conventional FL workflows.
\subsection{Economic Feasibility}
DAVS is economically feasible because it reduces recurring costs related to communication-heavy training and to failures caused by poisoned updates. Random Projection-based sketching reduces bandwidth per client per round, accelerating experimentation cycles and lowering network overhead. Compared to cryptography-heavy defenses, DAVS introduces low additional computation while providing robust filtering against malicious updates.

The permissioned blockchain audit layer stores lightweight metadata and hashes, keeping storage costs manageable. In the long term, preventing compromised models and enabling verifiable provenance can reduce financial and reputational risk associated with unsafe clinical AI deployments and compliance shortcomings.


\section{Project Management}
\subsection{Timeline Chart}
Below are the key milestones for the DAVS project, aligned with the design, implementation, and evaluation pipeline described in this report.

\begin{table}[h!]
\centering
\begin{tabular}{|p{4cm}|p{7cm}|p{4cm}|}
\hline
\textbf{Milestone} & \textbf{Description} & \textbf{Estimated Completion Date} \\
\hline
Project Initiation \& Planning & Define scope, datasets, threat model, and evaluation plan & August, 2025 \\
\hline
System Design \& Architecture & Define FL pipeline, sketching, DAVS selection, consensus, and audit logging & September, 2025 \\
\hline
Federated Learning Baseline & Implement baseline FedAvg training and non-IID partitioning & November, 2025 \\
\hline
Gradient Sketching Layer & Implement Random Projection compression and validate similarity preservation & December, 2025 \\
\hline
DAVS Committee Selection & Implement amnesic scoring and committee formation logic & January, 2026 \\
\hline
PBFT Validation \& Audit Ledger & Integrate PBFT-inspired validation and permissioned blockchain logging & February, 2026 \\
\hline
Attack Evaluation \& Benchmarking & Test poisoning attacks and compare against baselines & March, 2026 \\
\hline
Result Analysis \& Reporting & Consolidate results, tables, and discussion & March, 2026 \\
\hline
Final Review \& Documentation & Complete documentation and finalize project report & April, 2026 \\
\hline
Conference Presentation & Write research paper and present at ICTIS conference. & April, 2026 \\
\hline
\end{tabular}
\caption{DAVS Project Milestones and Timeline}
\end{table}

\\ 


\section{Project Resources}

\subsection{Hardware Requirements}
DAVS hardware requirements are driven by federated training workloads, local client training, and the overhead of validation and audit logging. These requirements ensure efficient experimentation and reliable execution of the end-to-end pipeline.

1. Development and Coordinator/Validator Server Requirements
Processor: Intel Core i7 (or equivalent AMD Ryzen 7) – 8th Gen or higher

RAM: Minimum 16 GB (Recommended: 32 GB for handling training orchestration, logging, and evaluation)

Storage: 1 TB SSD for dataset access, model checkpoints, and audit logs

GPU (Recommended): NVIDIA GPU (8--16 GB VRAM) for medical imaging model training

Internet: High speed broadband with upload/download speeds of at least 100 Mbps

Power Backup: UPS or generator backup to ensure continuous uptime during outages

2. Client Side (Hospital Training Nodes)
Processor: Intel Core i3 or higher

RAM: Minimum 4 GB (Recommended: 8 GB for smoother performance)

Storage: 250 GB HDD or SSD

Compatibility with Devices: laptops/desktops/servers capable of local training

Display: Minimum resolution of 1366 x 768 pixels

Internet: Stable connection (10 Mbps or higher) for transmitting sketches and receiving global models



\subsection{Software Requirements}

1. Operating System\\
Development Environment:
Windows 10/11 (best used when setting up blockchain nodes), macOS (use cross platform testing)
Client Environment (User Side):
Windows 10/11, Android 8.0+, iOS 13+, or any other operating system with a current web browser.

2. Frontend Technologies\\
Language: HTML5, CSS3, JavaScript
Framework: (Optional) lightweight monitoring interface for rounds and metrics.

3. Training \& Orchestration Stack\\
Runtime Environment: Python 3.x with ML libraries.
Frameworks/Libraries: PyTorch, NumPy; utilities for FL orchestration and evaluation.

4. Blockchain Audit Layer\\
Platform: Permissioned blockchain / custom ledger implementation for immutable round logging.

5. Security \& Validation\\
Mechanisms: Gradient sketching (Random Projection), DAVS committee selection, PBFT-inspired consensus, hashing for model provenance.

6. Development Deployment Tools\\
Version Control: Git with GitHub integration.
Package Managers: npm to deal with project dependencies.
IDE or Text Editors: Visual Studio Code or any development environment of choice.



\section{Project Estimation}

\subsection{COCOMO Estimation Model}
In order to determine the development effort of DAVS, the Basic COCOMO Model is used under Semi Detached category since the project has moderate complexity with the integration of ML training, consensus validation, and blockchain-based logging.
\\
\\
Estimated Size: 14,000 lines of code (14 KLOC)\\
Effort (E): $E = 3.0 \times (14)^{1.12} \approx 57 \text{ person-months}$\\
Development Time (D): $D = 2.5 \times (57)^{0.35} \approx 9 \text{ months}$\\
Team Size (P): $P = \frac{57}{9} \approx 6 \text{ members}$

\begin{table}[h!]
\centering
\begin{tabular}{|l|l|}
\hline
\textbf{Parameter}        & \textbf{Value}       \\
\hline
LOC                      & 14,000             \\
Effort                   & \textasciitilde57 PM    \\
Development Time         & \textasciitilde9 months \\
Team Size                & \textasciitilde6 members \\
\hline
\end{tabular}
\caption{Project Parameters}
\label{tab:project_parameters}
\end{table}


\subsection{Function Point Analysis}
Function Point Analysis (FPA) estimates the size and effort of DAVS by evaluating functional components such as client training, sketch generation, committee selection, consensus validation, audit logging, and metric reporting. The method does not depend on the programming language and is oriented towards system functionality.

1. Component Breakdown

\begin{table}[h!]
\centering
\begin{tabular}{|l|c|l|c|c|}
\hline
\textbf{Component}                                   & \textbf{Count} & \textbf{Complexity} & \textbf{Weight} & \textbf{Total FP} \\
\hline
External Inputs (client join, config, round control)              & 6     & Medium     & 4      & 24       \\
External Outputs (metrics, model hash, audit records)          & 4     & Medium     & 5      & 20       \\
User Inquiries (round status, committee logs)             & 3     & Simple     & 3      & 9       \\
Internal Logical Files (model states, logs, scores) & 4     & Medium     & 7      & 28       \\
External Interface Files (datasets, client nodes)             & 3     & Medium     & 7      & 21       \\
\hline
\textbf{Unadjusted Function Points}              &       &            &        & \textbf{102}  \\
\hline
\end{tabular}
\caption{Function Points Calculation}
\label{tab:function_points}
\end{table}

2. Complexity Adjustment
Based on the system features like distributed training, robustness validation, consensus, and audit integration, a Complexity Adjustment Factor (CAF) of 1.08 is used.
\\
Adjusted Function Points (AFP): $AFP = 102 \times 1.08 \approx 110$\\

3. Effort Estimation
With an average of 10 hours per function point:

Total Effort = 110 × 10 = 1,100 hours\\

A full-time (90 hours/month/person) 3-person team:

Time: $T = \frac{1100}{3 \times 90} \approx 4 \text{ months}$

\rule{0.49\textwidth}{0.75pt} $_{\bigcirc}$ \rule{0.49\textwidth}{0.75pt}\\
